% =====================================================================
%  chunks/ch03e.tex —— 原书 p.70–78（PDF 77–85）
%  第 3 章 §5 非单连通空间、模复形、Reidemeister 挠、简单同伦型
% =====================================================================
% [原文 p.70]
\begin{example}{2}
设 $\cF_1\subset\cF_2$ 是 $K$ 上的阿贝尔群常层，$\cF_3=\cF_1/\cF_2$。在这种情形下，(4.5) 给出 $K$ 的以不同系数群为系数的上同调群之间的自然关系。例如，若 $\cF_1$ 的常群为 $\Z/p$，$\cF_2$ 的为 $\Z/p^2$，则 $\cF_3$ 的为 $\Z/p^2/\Z/p\cong\Z/p$，而正合列 (4.5) 中的同态 $\partial^*$ 取形式
% [Tutor Note: 原书此例印作 $\cF_3=\cF_1/\cF_2$，但依上下文（$\cF_1\subset\cF_2$ 为常层，商层应为 $\cF_2/\cF_1$，且下文 $\Z/p^2/\Z/p$ 也按 $\cF_2/\cF_1$ 计算）此处应为 $\cF_2/\cF_1$；正文照排，不擅改。]
\begin{equation}
  \partial^*\colon H^q(K;\Z/p)\longrightarrow H^{q+1}(K;\Z/p). \tag{4.6}
\end{equation}
这个同态通常记为 $\beta$，称为 \emph{Bockstein 同态}。它的直接定义如下：设 $\tilde z$ 是 $C^*(K;\Z)$ 中的一个整 $q$-上闭链，它模 $p$ 约化后给出上闭链 $z\in C^*(K;\Z/p)$，即 $z=\tilde z\bmod p$。取上边缘得
\[
  \partial^*z=0,\qquad \partial^*\tilde z=pu,\qquad \partial^*u=0,
\]
于是表达式
\[
  \frac{1}{p}\partial^*\tilde z=\beta(z)
\]
对 $z\in H^q(K;\Z/p)$ 定义了 $H^{q+1}(K;\Z/p)$ 中的一个元素 $\beta(z)$。
\end{example}

\begin{remark}{}
值得指出，非阿贝尔群的层相当常见。对它们可以自然地、且有用处地定义带群结构的对象 $H^0(X;\cF)$，以及仅仅作为集合的 $H^1(X;\cF)$。对常层而言，这些对象实际上前面已经引入过（见 §2）。一般情形下的定义与此完全类似。下文讨论主 $G$-丛时我们会给出具体例子。群 $H^0(X;\cF)$ 称为层 $\cF$ 的\emph{截面群}，而 $\cF_U$ 则是区域 $U$ 上的截面（或截面芽）。
\end{remark}

\section{非单连通空间的同调论。模复形。Reidemeister 挠。简单同伦型}
\takeaway{非单连通复形的（上）同调带有独特而孤立的代数特征：把系数提升为 $\Z[\pi]$ 后，复形成为模复形，同调模被基本群作用所扭曲，由此产生挠与简单同伦型两类不变量。本节引入模复形、Reidemeister 挠与 Whitehead 群，说明它们如何介入同伦分类。}

非单连通复形的同调论与上同调论呈现出一些独特而孤立的代数特征，并由此在某些情形下引出重要结论。前面已经提到，对非单连通复形 $K$，可以给出对 $(K,L)$ 的 Hurewicz 定理的相对形式，它涉及模 $\pi_n(K,L)$ 在 $\Z[\pi]$ 上的结构，其中 $\pi=\pi_1(L)$。若把 $\Z$ 看作平凡 $\Z[\pi]$-模，即认为 $\pi$ 的元素都恒等地作用在 $\Z$ 上，则这一结论可以看作断言同构（对照 (3.18)）
\begin{equation}
  H_n(K,L;\Z)\cong\pi_n(K,L)\otimes_{\Z[\pi]}\Z, \tag{5.1}
\end{equation}
% [原文 p.71]
其成立条件为 $\pi_j(K,L)=0$（$0<j<n$）且 $n>1$。

给定环 $R$ 上的两个模 $M$、$N$，我们定义它们在 $R$ 上的\emph{张量积} $M\otimes_R N$ 为由所有符号 $m\otimes n$（$m\in M$，$n\in N$）生成、并满足通常双线性关系的模（对照 (2.3)）：
\begin{equation}
  \lambda(m\otimes n)=(\lambda m)\otimes n=m\otimes(\lambda n),\qquad \lambda\in R. \tag{5.2}
\end{equation}
当 $R=\Z$ 时，这就退化为 §2 中引入的阿贝尔群的张量积（把阿贝尔群自然看作 $\Z$-模）。

对任何非单连通的单纯复形（或 CW 复形）$K$，都自然联系着一个代数复形 $C_*(\tilde K)$，其中 $\tilde K$ 是 $K$ 的万有覆叠；它是自由 $\Z[\pi]$-模构成的复形，其基由与 $K$ 的单形（或胞腔）自然一一对应的链组成。复形 $C_*(\tilde K)$ 的定义如下。设
\[
  p\colon\tilde K\longrightarrow K,\qquad \bigl(\pi_1(\tilde K)=0\bigr),
\]
为万有覆叠投影。前面已经看到，群 $\pi=\pi_1(K)$ 自由地作用在 $\tilde K$ 上，而且把单形送到单形（把胞腔送到胞腔）；这一点来自如下事实：在可缩子空间之上，万有覆叠是平凡的，即该子空间与离散纤维 $F$ 的乘积，而这里的 $F$ 正是群 $\pi=\pi_1(K)$。在当前情形中，可缩子空间就是 $K$ 的单形（或胞腔），于是
\[
  p^{-1}(\sigma^n_\alpha)=\bigcup_{\gamma\in\pi}\sigma^n_{\alpha\gamma}.
\]
取一个「初始」$n$-单形（或胞腔）$\sigma^n_{\alpha 1}\subset p^{-1}(\sigma^n_\alpha)$。$p^{-1}(\sigma^n_\alpha)$ 在 $\tilde K$ 中的其余 $n$-单形，由群 $\pi$ 的元素作用在这一个上得到：
\[
  \gamma(\sigma^n_{\alpha 1})=\sigma^n_{\alpha\gamma},\qquad \gamma\in\pi.
\]
于是 $C_n(\tilde K)$ 就是 $\tilde K$ 的通常 $n$-链复形，并带有这个 $\Z[\pi]$-作用。它显然是自由模。注意 $\Z[\pi]$ 的作用与边缘算子 $\partial$ 交换：
\[
  \partial\lambda(c)=\lambda\partial(c),\qquad \lambda\in\Z[\pi].
\]
于是得到作为自由 $\Z[\pi]$-模的链复形 $C_*(\tilde K)$。

若覆叠 $\tilde K$ 不必是万有的，但仍为正规覆叠，则商群 $\Gamma=\pi_1(K)/N$（其中 $N\cong\pi_1(\tilde K)$）自由地作用在 $\tilde K$ 上，并把单形（胞腔）送到单形（胞腔），使轨道空间 $\tilde K/\Gamma\cong K$。此时与前一样，链复形 $C_*(\tilde K)$ 带有自然的 $\Gamma$-作用，从而成为自由 $\Z[\Gamma]$-模。

% [原文 p.72]
群 $\Gamma$ 的任一表示 $\rho$，即任一从 $\Gamma$ 到某个向量空间或阿贝尔群 $V$ 的自同构群 $\operatorname{Aut}V$ 的同态 $\rho\colon\Gamma\to\operatorname{Aut}V$，都可用来给 $V$ 赋予 $\Z[\Gamma]$-模的结构 $(V,\rho)$。把这样的模 $(V,\rho)$ 与刚才定义的 $\Z[\Gamma]$-模 $C_*(\tilde K)$ 作张量积，就得到 $K$ 关于 $\rho$ 的\emph{链复形}
\begin{equation}
  C_*(K;\rho)=C_*(\tilde K)\otimes_{\Z[\Gamma]}(V,\rho). \tag{5.3}
\end{equation}
由此按相应的商模构造出 $K$ 关于表示 $\rho$ 的同调群与上同调群。

\begin{example}{1}
若取 $\rho$ 为把 $\Gamma$ 中所有元素都送到 $V$ 的恒等自同构的平凡表示，则得到通常的同调群与上同调群 $H_q(K;V)$、$H^q(K;V)$。
\end{example}

\begin{example}{2}
设 $K$ 是不可定向流形，$\dim K=n$。对每个道路类 $[\gamma]\in\pi_1(K)$，定义 $\rho[\gamma]=\sgn\gamma$，其中若沿 $\gamma$ 搬运坐标标架保持定向则 $\sgn\gamma=1$，否则为 $-1$。取 $V=\Z$，此时 $\operatorname{Aut}\Z\cong\Z/2\cong\{\pm1\}$。所得同调群与上同调群 $H_j(K;\rho)$、$H^j(K;\rho)$ 之间，存在如前面由 Poincaré 对偶复形 $DK$ 所确定的 Poincaré 对偶同构（见 §2）：
\begin{equation}
\begin{aligned}
  D\colon&\ H_j(K;\rho)\cong H^{n-j}(K;\rho),\\
  D\colon&\ H^j(K;\rho)\cong H_{n-j}(K;\rho).
\end{aligned}\tag{5.4}
\end{equation}
\end{example}

\begin{example}{3}
取 $V=\Z[\pi]$，并让 $\pi$ 的每个元素 $g$ 通过右乘作用在 $\Z[\pi]$ 上。这给出实质上就是恒等的表示 $\rho$，从而
\[
  H_j(K;\rho)\cong H_j(\tilde K;\Z),\qquad H^j(K;\rho)\cong H^j(\tilde K;\Z).
\]
\end{example}

特别令人感兴趣的是那些表示 $\rho\colon\pi\to\operatorname{Aut}V$，其相应的同调群 $H_q(K;\rho)$ 都平凡；此时称 $K$ 为 \emph{$\rho$-零调的}复形。若 $K$ 是 $n$ 维的 $\rho$-零调复形，则链复形 $C_*(K;\rho)$ 给出正合列
\begin{equation}
  0\longrightarrow C_n(K;\rho)\xrightarrow{\partial}C_{n-1}(K;\rho)\xrightarrow{\partial}\cdots\xrightarrow{\partial}C_1(K;\rho)\xrightarrow{\partial}C_0(K;\rho)\longrightarrow 0. \tag{5.5}
\end{equation}

下文中重要的是：若 $V$ 是带给定自由基的自由阿贝尔群（或带给定基的向量空间），则自由阿贝尔群 $C_j(K;\rho)$（当 $V$ 是向量空间时为向量空间）也各自带有\emph{特异基}，其成员与 $K$ 的胞腔相对应。

\begin{remark}{}
设 $V$ 是带给定基的自由阿贝尔群（或向量空间）。刚才描述的 $C_j(K;\rho)$ 的「特异基」，在至多相差下列初等基变换之一的意义下唯一确定：

% [原文 p.73]
1) 由某些或全部胞腔的定向改变引起的变换：$\sigma^n_\alpha\mapsto-\sigma^n_\alpha$；

2) 由从完全原像 $p^{-1}(\sigma^n_\alpha)=\bigcup\sigma^n_{\alpha\gamma}$（$\gamma\in\pi$）中另选「初始胞腔」$\sigma^n_{\alpha 1}$ 引起的变换；此时基变换的形状为 $\sigma^n_\alpha\mapsto\pm\gamma\sigma^n_\alpha$，或者相对于表示 $\rho$ 为 $\sigma^n_\alpha\mapsto\pm\rho(\gamma)\sigma^n_\alpha$。

以下我们假定对每个 $C_j(K;\rho)$ 都固定一个「特异基」。\hfill$\square$
\end{remark}

现在设 $V$ 是某个域 $k$ 上的有限维向量空间，带有给定的基；又设 $K$ 是 $\rho$-零调的，于是有正合列 (5.5)。若该正合列只有两项
\[
  0\longrightarrow C_n(K;\rho)\xrightarrow{\partial}C_{n-1}(K;\rho)\longrightarrow 0,
\]
则 $\partial$ 就是同构，相对于 $C_n(K;\rho)$ 与 $C_{n-1}(K;\rho)$ 的特异基，它由 $k$ 上的非奇异矩阵表示，且 $\det\partial=\det C_*(K;\rho)$。可以证明，对每个有限的 $\rho$-零调复形 $K$，都可以相对于相应正合列 (5.5) 中各个 $C_j(K;\rho)$ 的特异基，联系一个「行列式」。当正合列 (5.5) 有三项时，
\[
  0\longrightarrow C_n(K;\rho)\xrightarrow{\partial}C_{n-1}(K;\rho)\xrightarrow{\partial}C_{n-2}(K;\rho)\longrightarrow 0,
\]
这个行列式定义如下。把 $C_n(K;\rho)$、$C_{n-1}(K;\rho)$、$C_{n-2}(K;\rho)$ 的特异基分别记为 $c_n$、$c_{n-1}$、$c_{n-2}$。中间的复形 $C_{n-1}(K;\rho)$ 还有另一组基 $\bigl(\partial c_n,\partial^{-1}(c_{n-2})\bigr)$，其中 $\partial^{-1}(c_{n-2})$ 表示 $c_{n-2}$ 的成员的一组任意选定的原像，并按对应方式排序。由于选择 $\partial^{-1}(c_{n-2})$ 的自由度具有「三角形」性质，从 $c_{n-1}$ 到 $\bigl(\partial c_n,\partial^{-1}(c_{n-2})\bigr)$ 的基变换行列式是唯一确定的。一般情形下，行列式 $\det C_*(K;\rho)$ 类似地定义：把每个 $C_j=C_j(K;\rho)$ 的特异基记为 $c_j$，用 $A_j\subset C_{j-1}$ 表示 $C_j$ 在 $\partial$ 下的像（$j=1,\dots,n-1$），并为这些 $A_j$ 任意选取基 $a_j$。令 $A_0=0$、$A_n=C_n$、$a_n=c_n$：
\[
  A_n=c_n,\qquad A_j=\partial C_j\subset C_{j-1},\quad j=1,\dots,n-1,\qquad A_0=0,
\]
\[
  0\longrightarrow A_n\longrightarrow C_{n-1}\xrightarrow{\partial}A_{n-1}\subset C_{n-2},
\]
\[
  0\longrightarrow A_{n-1}\longrightarrow C_{n-2}\xrightarrow{\partial}A_{n-2}\subset C_{n-3},
\]
\[
  \cdots\cdots\cdots\cdots\cdots\cdots\cdots\cdots\cdots\cdots
\]
\[
  0\longrightarrow A_1\longrightarrow C_0\xrightarrow{\partial}A_0=0.
\]

% [原文 p.74]
我们于是定义
\[
  \det C_*(K;\rho)=
  \frac{\det[c_{n-1}:(a_n,\partial^{-1}a_{n-1})]\cdot\det[c_{n-3}:(a_{n-2},\partial^{-1}a_{n-3})]\cdots}
       {\det[c_{n-2}:(a_{n-1},\partial^{-1}a_{n-2})]\cdot\det[c_{n-4}:(a_{n-3},\partial^{-1}a_{n-4})]\cdots},
\]
这里 $C_*(K;\rho)$ 是任意的 $\rho$-零调复形。

\begin{definition*}{5.1}
上面为 $\rho$-零调复形 $K$（即满足 $H_j(K;\rho)=0$，$j\ge0$）所定义的量 $\det C_*(K;\rho)$，称为 $K$ 的 \emph{Reidemeister 挠}，记为 $R(K,\rho)$。
\end{definition*}

20 世纪 30 年代末（由 Reidemeister、de Rham、Franz）已经确立：在相差一个符号、并乘以任意多个形如 $\det\rho(\gamma)$（$\gamma\in\pi$）的数之后，量 $\det C_*(K;\rho)=R(K,\rho)$ 是单纯复形 $K$ 的组合不变量，即它在重分下不变。直到很久以后，20 世纪 70 年代才证明 $R(K,\rho)$ 实际上是拓扑不变量。（$R(K,\rho)$ 的这种变动，来自前面提到的选取 $C_j$ 特异基时的自由度。）

% [Tutor Note: 注意挠不是数，而是「在乘以 $\pm\det\rho(\gamma)$ 的意义下」确定的量：这些因子来自基的选取（胞腔定向、初始胞腔的选取以及表示 $\rho$）。谈挠是否相等时，必须先把这个歧义固定下来；这也解释了为何挠的早期结果只能给出相差一个乘性因子的分类。]

还须注意，若 $L$ 是任一单连通复形，$\pi_1L=0$，则有
\[
  R(K\times L,\rho)=R(K,\rho)\chi(L),
\]
其中 $\chi(L)$ 是 $L$ 的 Euler--Poincaré 示性数。

量 $R(K,\rho)$ 最初（由 Reidemeister）引入，目的是研究「透镜空间」，后者定义如下。设 $A$ 是 $n\times n$ 复对角矩阵，其对角元为本原 $m$ 次单位根：
\begin{equation}
  A=(a_{jk}),\qquad a_{jk}=\ee^{2\pi\ii q_j/m}\delta_{jk},\qquad A^m=I_{n\times n}, \tag{5.6}
\end{equation}
其中 $(\delta_{jk})=I_{n\times n}$ 是 $n\times n$ 单位矩阵。由于 $a_{jj}$ 是本原 $m$ 次单位根，每个 $q_j$ 都必须与 $m$ 互素。矩阵 $A$ 生成一个同构于 $\Z/m$ 的群，它按如下规则作用在球面 $S^{2n-1}$ 上：
\begin{equation}
  z^j\xrightarrow{A}\ee^{2\pi\ii q_j/m}z^j,\qquad \sum_{j=1}^{n}|z^j|^2=1. \tag{5.7}
\end{equation}
取轨道空间，就得到称为 $(2n-1)$ 维\emph{透镜空间}的流形：
\begin{equation}
  L^{2n-1}_{(q_1,\dots,q_n)}=S^{2n-1}/A. \tag{5.8}
\end{equation}
总可以假定 $q_1=1$：只需用 $A$ 所生成循环群的适当生成元代替 $A$。

当 $n=2$ 时，得到 3 维流形 $L^3_q=L^3_{(1,q)}$，其中 $q$ 与 $m$ 互素。两个这样的流形 $L^3_{q'}$ 与 $L^3_{q''}$（对同一个 $m$ 定义）同伦等价，当且仅当
\begin{equation}
  q'/q''\equiv\pm\lambda^2\ \operatorname{mod} m. \tag{5.9}
\end{equation}

% [原文 p.75]
这一点可借助上同调环 $H^*(L^3_q;\Z/m)$ 的结构以及 Bockstein 同态 $\beta\colon H^1(L^3_q;\Z/m)\to H^2(L^3_q;\Z/m)$（见 (4.6)）直接得到：对透镜空间 $L^3_q$ 的基本同调类 $[L^3_q]$，成立同余式
\[
  \bigl(a\beta(a),\pm[L^3_q]\bigr)\equiv\pm q\ \operatorname{mod} m,
\]
其中 $a$ 是 $H^1(L^3_q;\Z/m)\cong\Z/m$ 的任一生成元。若把这个同余式中的 $a$ 换成任一其他生成元 $\lambda a$，则 $q$ 随之变为 $\lambda^2 q$。

透镜空间 $L^{2n-1}_{(q_1,\dots,q_n)}$ 可以分解为 CW 复形，而且分解方式与 $(q_1,\dots,q_n)$ 无关，这与射影空间 $\RP^{2n-1}$ 的作法颇为相似。这样每个维数上恰有一个胞腔：$\sigma^0,\sigma^1,\dots,\sigma^{2n-1}$，其边缘算子由下式给出
\begin{equation}
  \partial\sigma^{2j-1}=0,\qquad \partial\sigma^{2j}=m\sigma^{2j-1},\qquad j>0. \tag{5.10}
\end{equation}

然而在球面 $S^{2n-1}$ 上，由 $A$ 生成的 $m$ 阶循环群自由且保胞腔地作用；作为覆叠空间，它在 $S^{2n-1}$ 上诱导的胞腔分解确实本质地依赖于 $(q_1,\dots,q_n)$。当 $n=2$ 时，由 $L^3_q$ 的胞腔在 $S^{2n-1}=S^3$ 上诱导出如下胞腔：
\begin{equation}
  \sigma^s_j=A^j\sigma^s,\qquad s=0,1,2,3,\quad j=0,1,\dots,m-1. \tag{5.11}
\end{equation}
边缘算子在这些胞腔上的作用如下（记 $\sigma^s_0=\sigma^s$，$s=0,1,2,3$）：
\begin{equation}
\begin{aligned}
  \partial\sigma^3&=(1-A^q)\sigma^2, & \sigma^3_1&=\sigma^3,\\
  \partial\sigma^2&=\bigl(1+A+\cdots+A^{m-1}\bigr)\sigma^1,\\
  \partial\sigma^1&=(1-A)\sigma^0.
\end{aligned}\tag{5.12}
\end{equation}

对表示 $\rho\colon\Z\to(\ee^{2\pi\ii/m})$（在空间 $V=\C^1$ 中），透镜空间 $L^3_q$ 是 $\rho$-零调的，由复形 $C_*(L^3_q;\rho)$ 得到的正合列 (5.5) 为
\[
  0\longrightarrow C_3\xrightarrow{\partial}C_2\xrightarrow{\partial}C_1\xrightarrow{\partial}C_0\longrightarrow 0,
\]
其中每个 $C_j$ 都是一维的，相应地有特异生成元，记为 $e_j$。由 (5.12) 知边缘算子为
\begin{equation}
\begin{aligned}
  \partial e_3&=\bigl(1-\ee^{2\pi\ii q/m}\bigr)e_2, & \partial e_2&=0,\\
  \partial e_1&=\bigl(1-\ee^{2\pi\ii q/m}\bigr)e_0,
\end{aligned}\tag{5.13}
\end{equation}
由此可算出 Reidemeister 挠，得
\begin{equation}
  R(L^3_q,\rho)=(1-\zeta^q)(1-\zeta)\ \operatorname{mod}\ (\pm\zeta^\alpha), \tag{5.14}
\end{equation}

% [原文 p.76]
其中 $\zeta=\ee^{2\pi\ii/m}$。由 (5.14) 可推出：存在同伦等价、但不同伦组合等价的透镜空间。其后（Moise、Chapman 与 Edwards 的）结果表明，同伦等价的透镜空间事实上甚至不必同胚。de Rham 与 Franz 通过考虑 $n>2$ 时的各种表示 $\rho$ 证明：事实上只有在平凡的情形下，透镜空间才组合等价。

对这一类问题的完整分析，是 Whitehead 从同伦论观点作出的。设已知两个有限复形 $K$ 与 $L$ 同伦等价。能否借助初等的组合运算，把这一同伦等价实现出来？我们不去对这类运算作详细而精确的几何描述，而只限于勾画其本质的代数图景。

设 $L$ 作为单纯子复形或 CW 子复形含于 $K$ 中。由于 $K\sim L$，作为 $\Z[\pi]$-模看待的相对链复形 $C(K,L)=\bigoplus_t C_t(K,L)$ 是零调的。考察任意结合环 $R$ 上任意代数复形 $C=\bigoplus_t C_t$ 上的如下两类初等运算，其中每个 $C_t$ 都带特异基 $\{\sigma^t_i\}$：

(i) 对每个固定的 $k$、$j$ 与 $\gamma\in R$，
\begin{equation}
  \sigma^k_j\mapsto\gamma\sigma^k_j+\sigma^k_j,\qquad
  \sigma^t_i\mapsto\sigma^t_i\quad\text{若 } t\neq k \text{ 或 } i\neq j. \tag{5.15}
\end{equation}

(ii) 稳定化：对每个固定的 $k$，给 $C$ 添上一个形如下的直和项 $\tilde C^{(k)}$：
\begin{equation}
  0\longrightarrow\tilde C^{(k)}_k\xrightarrow{\partial}\tilde C^{(k)}_{k-1}\longrightarrow 0, \tag{5.16}
\end{equation}
其中 $\tilde C^{(k)}_{k-1}$ 与 $\tilde C^{(k)}_k$ 是循环的，分别由 $\tilde e_{k-1}$、$\tilde e_k$ 生成，且 $\partial\tilde e_k=\tilde e_{k-1}$；于是该运算定义如下：
\[
  C\longrightarrow C\oplus\tilde C^{(k)},
\]
即
\begin{equation}
  C_k\longrightarrow C_k\oplus\tilde C^{(k)}_k,\qquad C_{k-1}\longrightarrow C_{k-1}\oplus\tilde C^{(k)}_{k-1},\qquad C_j\longrightarrow C_j\ \ (j\neq k-1,k). \tag{5.17}
\end{equation}

若定义 $R$ 上两个零调代数复形\emph{等价}，是指其中一个可由另一个经有限次运算 (i)、(ii) 及其逆运算得到，则在取直和下，所得的等价类构成一个阿贝尔群 $K_1(R)$。这就是「行列式群」；它也可以等价地定义为 $R$ 上无穷矩阵群的某个商群，这些矩阵只有有限多个非零的非对角元，且除有限多个对角元外，其余对角元都等于 $R$ 的乘法单位元。为得到 $K_1(R)$，须把这个群对由所有换位子以及所有如下的初等矩阵所生成的子群取商：
% [原文 p.77]
这些矩阵的对角元为 1，唯一的其余非零元位于对角线上方。这个群在代数中自然地出现，它是把行列式的基本性质公理化之后的产物。

在当前拓扑的语境中，环 $R$ 形如 $\Z[\pi]$；考虑到前面所说的选取特异基时的自由度——即由覆叠空间胞腔的定向以及初始胞腔的选取所带来的自由度——重要的其实不是 $K_1(\Z[\pi])$，而是它的商群
\[
  \operatorname{Wh}(\pi)=K_1(\Z[\pi])/(\pm\pi)
\]
（即 \emph{Whitehead 群}）。

于是对同伦等价复形的对 $L\subset K$，我们得到不变量：$\Z[\pi]$-模复形 $C(\tilde K,\tilde L)$ 在 Whitehead 群 $\operatorname{Wh}(\pi)$ 中的「行列式」。与此相关的术语是「简单同伦型」：若其中一个复形能借助上面定义的初等组合运算变成另一个，则称这两个复形有相同的\emph{简单同伦型}。任一表示 $\rho\colon\pi\to\operatorname{GL}(n,\R)$ 都给出从 $K_1(\Z[\pi])$ 到 $\R$ 的乘法群的同态
\[
  K_1(\Z[\pi])\longrightarrow\R^*,
\]
由此得到前面讨论的 Reidemeister 挠。对单连通复形，群 $K_1$ 是平凡的：$\Z[\pi]=\Z$，$K_1(\Z)=0$，简单同伦等价就退化为通常的同伦等价。

早在 20 世纪 30 年代末，从 Smith 与 Richardson 的工作开始，人们就在研究有限变换群的同调性质。设有限群 $G$ 作用在 CW 复形 $K$ 上，并把胞腔送到胞腔。由这一作用，与前面颇为类似地得到一个 $\Z[G]$-模复形——即使 $G$ 的作用并不自由。把 $K$ 中由 $G$ 的每个元素都不动的点构成的子复形记为 $K^G$，相应的轨道复形记为 $K/G$。我们有自然映射 $i\colon K^G\to K$、$\pi\colon K\to K/G$，以及诱导同态
\[
  H_*(K^G)\xrightarrow{i_*}H_*(K),\qquad H_*(K)\xrightarrow{\pi_*}H_*(K/G).
\]

$G$ 的有限性使人可以定义\emph{转移同态}
\[
  \mu_*\colon H_*(K/G)\longrightarrow H_*(K),
\]
（它对任意系数群都有定义）如下：$C_j(K/G)$ 中的每个 $j$-链 $c_G$ 对应 $C_j(K)$ 中 $j$-链在有限群 $G$ 作用下的一个轨道 $\{g(c)\mid g\in G\}$，其中 $c\in C_j(K)$；同态 $\mu_j$ 则由如下对应所诱导：把每个 $c_G\in C_j(K/G)$ 对应到 $\sum_{g\in G}g(c)\in C_j(K)$，这里出现的和是相应轨道中各链的（有限）和。
% [原文 p.78]
由此可知，复合映射
\begin{equation}
  \pi_*\mu_*\colon H_*(K/G)\longrightarrow H_*(K/G) \tag{5.18}
\end{equation}
就是乘以 $|G|$；而
\begin{equation}
  \mu_*\pi_*\colon H_*(K)\longrightarrow H_*(K) \tag{5.19}
\end{equation}
由 $\sum_{g\in G}g$ 的作用决定（当系数群为 $\Z$ 时，该元素属于 $\Z[G]$）。由 (5.18) 与 (5.19) 推出：在特征为零、或有限特征与 $|G|$ 互素的域上，有
\begin{equation}
\begin{aligned}
  H_*(K;k)&\cong\mu_*H_*(K/G;k)\oplus\Ker\pi_*,\\
  H_*(K/G;k)&\cong H_*(K;k)^G,
\end{aligned}\tag{5.20}
\end{equation}
其中上标 $G$ 表示被 $G$ 逐点固定的子群。
% [Tutor Note: 转移同态 $\mu_*$ 的要点是 $G$ 有限：把轨道上的链求和才得到定义在 $K/G$ 上的链，而复合 $\pi_*\mu_*$ 恰是乘以 $|G|$。这也解释了对系数域的要求：若域特征整除 $|G|$，$\pi_*\mu_*$ 不再是同构，(5.20) 的分解随之失效。]

借助 $G$-作用，$\Z[G]$ 的每个元素 $\rho$ 都联系着一个链复形 $\rho C(K)$ 及相应的同调群 $H^\rho_*(K)=H_*(\rho C(K))$。情形 $G=\Z/p$、$k=\Z/p$ 特别令人感兴趣。注意这里 $H^\rho_*(K;\Z/p)$ 应理解为复形 $\rho C(K;\Z/p)$ 的同调，而不是 $\rho C(K;\Z)\otimes\Z/p$ 的同调。考察群环 $k[G]$（$k=\Z/p$）中的下列元素：
\begin{equation}
\begin{aligned}
  \sigma&=\sum_{g\in G}g=1+T+T^2+\cdots+T^{p-1},\\
  \tau&=1-T,\qquad \sigma\tau=\tau\sigma=0,
\end{aligned}\tag{5.21}
\end{equation}
其中 $T$ 是 $p$ 阶循环群 $G$ 的生成元（用乘法记号）。Smith 的正合列如下：
\begin{equation}
\begin{aligned}
  \cdots\longrightarrow H_j(K;\Z/p)&\xrightarrow{\sigma_*}H^\sigma_j(K;\Z/p)\xrightarrow{\partial}H^\tau_{j-1}(K;\Z/p)\oplus H_{j-1}(K^G;\Z/p)\xrightarrow{i_*}\\
  &H_{j-1}(K;\Z/p)\xrightarrow{\sigma_*}\cdots;
\end{aligned}\tag{5.22}
\end{equation}
\[
\begin{aligned}
  \cdots\longrightarrow H_j(K;\Z/p)&\xrightarrow{\tau_*}H^\tau_j(K;\Z/p)\xrightarrow{\partial}H^\sigma_{j-1}(K;\Z/p)\oplus H_{j-1}(K^G;\Z/p)\xrightarrow{i_*}\\
  &H_{j-1}(K;\Z/p)\xrightarrow{\tau_*}\cdots.
\end{aligned}
\]
存在类似的上同调正合列（箭头方向相反）。

由这些正合列可推出如下结果：若群 $\Z/p$ 作用在 $\Z/p$-同调的 $n$-球面上，则该作用的不动点集同样是某个维数 $r$ 的 $\Z/p$-同调球面。

% CHUNK ch03e DONE
